\documentclass[10pt,twocolumn]{article}
\usepackage[T1]{fontenc}
\usepackage[utf8]{inputenc}
\usepackage{lmodern}
\usepackage[margin=1.8cm,top=2.4cm,bottom=2cm,columnsep=0.6cm]{geometry}
\usepackage{fancyhdr}
\usepackage{titlesec}
\usepackage{booktabs}
\usepackage{amsmath,amssymb,amsfonts}
\usepackage{microtype}
\usepackage{xcolor}
\usepackage{mdframed}
\usepackage{parskip}
\usepackage{enumitem}
\usepackage{hyperref}

\definecolor{rulered}{RGB}{180,30,30}
\definecolor{boxgray}{RGB}{245,245,245}
\definecolor{keywordgray}{RGB}{100,100,100}

\pagestyle{fancy}
\fancyhf{}
\fancyhead[L]{\textsc{\small Claude Mag}}
\fancyhead[C]{\small\textit{Machine Learning Research Digest}}
\fancyhead[R]{\small 2026-07-20}
\fancyfoot[C]{\small\thepage}
\renewcommand{\headrulewidth}{0.8pt}

\titleformat{\section}{\normalsize\bfseries\color{rulered}}{}{0pt}{}%
  [\vspace{-4pt}\color{rulered}\rule{\columnwidth}{0.4pt}]
\titlespacing{\section}{0pt}{8pt}{4pt}

\hypersetup{colorlinks=true,urlcolor=rulered,linkcolor=black}

\begin{document}
\setlength{\parindent}{0pt}
\setlength{\parskip}{4pt}

{\small\color{keywordgray}\textbf{Keywords:}
  credit assignment \textbullet{}
  long-horizon agents \textbullet{}
  reinforcement learning \textbullet{}
  tool use}\par\vspace{4pt}

{\LARGE\bfseries\raggedright
  TRACE}\par\vspace{4pt}

{\large\itshape\raggedright
  Microsoft Research Derives Dense Turn-Level Rewards for
  Long-Horizon Agents from Temporal-Difference Changes in
  Frozen Reference Model Log-Probabilities}\par\vspace{6pt}

{\small\textbf{By} Leitian Tao, Baolin Peng, Wenlin Yao, Tao Ge,
  Hao Cheng, Mike Hang Wang, Jianfeng Gao, Sharon Li
  (Microsoft Research; University of Wisconsin-Madison)}\par\vspace{2pt}

{\small\color{keywordgray}arXiv:2607.13988 \textbullet{}
  July 15, 2026}
\par\vspace{2pt}\noindent{\color{rulered}\rule{\columnwidth}{0.6pt}}\vspace{6pt}

Multi-turn tool-use agents have emerged as a dominant paradigm for
complex task completion: the agent writes and executes code, calls
APIs, browses the web, edits files, and iterates over dozens to
hundreds of steps before producing a final answer. Outcome-only
reinforcement learning provides a single binary signal at the end of
this trajectory. TRACE (Turn-level Reward Assignment via Credit
Estimation) replaces this single sparse signal with dense,
turn-by-turn rewards derived from how each tool call changes the
model's probability of reaching the correct answer --- without any
human annotation of intermediate steps.

\begin{mdframed}[backgroundcolor=boxgray,linecolor=rulered,linewidth=1pt,
    innerleftmargin=6pt,innerrightmargin=6pt,
    innertopmargin=6pt,innerbottommargin=6pt]
\textbf{\small AT A GLANCE}\par\vspace{4pt}
\begin{description}[leftmargin=0pt,labelsep=4pt,itemsep=2pt,parsep=0pt,
    font=\small\bfseries]
  \item[Problem:] Outcome rewards become sparse, high-variance, and
    misleading for agents operating over 20--100 tool-call trajectories.
  \item[Why it matters:] Long-horizon agents are increasingly deployed
    for software engineering, scientific computing, and complex
    information retrieval; outcome-only RL trains them poorly.
  \item[Method:] State values estimated as log-ratio of gold-answer
    log-probabilities (current vs.\ frozen reference); turn-level
    rewards are TD differences between consecutive state values.
  \item[Result:] ToolBench task completion: 41.3\%~$\to$~58.7\%;
    WebArena: 19.4\%~$\to$~28.1\%; long-horizon code repair: 31.6\%~$\to$~47.2\%.
  \item[Evidence:] ToolBench, WebArena, and a 20-step code repair
    benchmark; no intermediate human annotations used.
\end{description}
\end{mdframed}

\section{The Big Picture}

Reinforcement learning for single-turn tasks is straightforward:
an answer is either right or wrong, and the reward is binary.
For multi-turn agents, the binary outcome arrives only at the
end of a trajectory that may span dozens of tool calls and
intermediate reasoning steps. Any single turn --- a web search,
a code edit, a database query --- affects the probability of
the final correct answer, but outcome-only RL provides no
per-turn signal about which turns were helpful.

The credit assignment problem is further complicated by
the misleading nature of failed trajectories: a sequence of
mostly correct intermediate steps followed by a single error
at the last step receives the same negative reward as a
trajectory that was wrong from the start, even though the
former's intermediate actions should be reinforced.

\section{Why Existing Approaches Fall Short}

Monte Carlo rollout-based credit assigns turn-level rewards
by averaging outcomes over multiple rollouts starting from
each intermediate state. This is statistically sound but
computationally prohibitive for long trajectories: estimating
a reliable MC signal for each of 50 turns in a trajectory
requires $50 \times K$ rollouts per training example.

Process reward models (PRMs) learn to score intermediate
steps, but require human annotation of intermediate step
quality, which is expensive and does not scale to the diverse
tool-use environments where agents are deployed.

Hindsight relabeling and advantage-based credit methods
improve over uniform assignment but remain high-variance for
very long trajectories because the value function must
extrapolate many steps into the future.

\section{Core Mechanism}

TRACE introduces a principled value function for multi-turn
trajectories that requires no additional training: the
log-ratio of the gold-answer log-probability under the
current policy versus a frozen reference model.

At each state $s_t$ (the trajectory prefix through turn $t$),
the state value is:
\[
V(s_t) = \log \frac{p_\theta(y^* \mid s_t)}{p_{\text{ref}}(y^* \mid s_t)}
\]
where $y^*$ is the known gold answer. This quantity has a natural
interpretation: it measures how much more likely the current
policy is (relative to the frozen reference) to produce the
correct answer given the information available at turn $t$.

The turn-level reward is the Temporal-Difference change:
\[
r_t = V(s_t) - V(s_{t-1})
\]

\section{Technical Formulation}

A trajectory of $T$ turns decomposes into state transitions:
$(s_0, a_1, s_1, a_2, \ldots, a_T, s_T)$, where each action $a_t$
is a tool call and $s_t$ is the trajectory prefix after the tool
result. The policy $\pi_\theta$ maps from $s_{t-1}$ to $a_t$; the
tool environment is deterministic.

The TRACE reward satisfies a telescoping property:
\[
\sum_{t=1}^{T} r_t = V(s_T) - V(s_0)
\]
where $V(s_T)$ approximates the outcome reward and $V(s_0)$ is a
constant baseline. Thus TRACE distributes the outcome signal across
turns in a way that is globally consistent: turns that increase
$V$ receive positive rewards, turns that decrease it receive negative
rewards, and the sum equals the net outcome advantage.

Log-probabilities $p_\theta(y^* \mid s_t)$ are computed by
teacher-forcing the gold answer given the trajectory prefix, a
standard operation for autoregressive models that requires one
forward pass per turn.

\section{Experimental Findings}

\begin{center}
\small
\begin{tabular}{lcc}
\toprule
Benchmark & Outcome RL & TRACE RL \\
\midrule
ToolBench & 41.3\% & \textbf{58.7\%} \\
WebArena & 19.4\% & \textbf{28.1\%} \\
Code Repair ($>$20 steps) & 31.6\% & \textbf{47.2\%} \\
\bottomrule
\end{tabular}
\end{center}

Credit quality measured against human annotations (Pearson $r$):
TRACE achieves $r = 0.61$ vs.\ $r = 0.22$ for MC rollout-based
credit and $r = 0.09$ for uniform assignment.

\section{Ablations and Interpretation}

\textbf{Uniform credit (assigning $V(s_T)/T$ to each turn):}
Performance degrades by 8--12\% across benchmarks, confirming
that the temporal structure of the TD signal is informative ---
not just the magnitude of the outcome.

\textbf{Reference model choice:} Using a weaker (smaller) reference
model degrades credit quality, suggesting the value of a calibrated
log-probability baseline. Using the same model checkpoint before RL
(as TRACE prescribes) consistently outperforms alternatives.

\textbf{Trajectory length:} The TRACE advantage grows with
trajectory length. For tasks with fewer than 5 tool calls,
the improvement over outcome-only RL is $<$2 percentage points.
For tasks with 20+ tool calls, the gap is 15+ percentage points,
consistent with the hypothesis that credit assignment difficulty
scales with trajectory length.

\textbf{Negative turn identification:} TRACE correctly assigns
negative rewards to turns that move the agent away from the goal
--- detectable as a decrease in state value --- enabling the
policy to specifically learn to avoid those actions.

\section*{Reference}

Leitian Tao et al.\
\textbf{TRACE: Turn-level Reward Assignment via Credit Estimation
for Long-Horizon Agents}.
arXiv:2607.13988, July 2026.
\url{https://arxiv.org/abs/2607.13988}

\end{document}
